% !TeX root = ../概率论与数理统计笔记.tex
% 由主文档通过 \input 载入；请编译主文档。
\setcounter{chapter}{0} % 保留原书章号。
\chapter{随机试验、事件与概率空间}
\label{chap:ps-events}

本章从“怎样把一次随机试验写成数学对象”出发，建立概率论的基本语言。
学习目标是正确选择样本空间、表达事件，并说明所采用的概率模型为什么适合问题。先修知识为集合运算与初等计数。

\section{随机试验与事件的集合表示}
\label{sec:ps-events-language}
引入\emph{随机试验}（random experiment）、\emph{样本空间}（sample space）与\emph{事件}（event）。
以抛硬币、抽球和等待时间说明有限、可数与不可数样本空间；讲解事件的并、交、补、差以及德摩根律。
区分互斥事件、对立事件、必然事件与不可能事件，并训练把“至少”“至多”“恰好”翻译为集合关系。

\section{古典概型与计数方法}
\label{sec:ps-counting}
在有限样本空间且基本结果等可能的条件下建立\emph{古典概型}（classical probability model）。
整理加法原理、乘法原理、排列、组合与有限容斥公式；比较有放回和无放回、有序和无序抽样。
先验证等可能性，再使用有利结果数与总结果数之比，避免对未经论证的结果分类直接等权计数。

\section{概率公理与基本性质}
\label{sec:ps-probability-axioms}
定义\emph{$\sigma$ 代数}（sigma-algebra）$\mathcal F$ 与\emph{概率空间}（probability space）$(\Omega,\mathcal F,\mathbb P)$，
说明非负性、规范性与可数可加性。
由公理证明补事件公式、单调性、有限加法公式、并集上界以及概率对单调事件列的连续性；强调一般样本空间中事件属于指定的事件族。

\section{几何概率与建模边界}
\label{sec:ps-geometric-probability}
在有限且非零的长度、面积或体积区域上，说明均匀分布假设如何导出\emph{几何概率}（geometric probability）。
比较不同随机抽取机制产生的不同模型，区分概率为零与不可能、概率为一与逻辑必然。
为连续型随机变量中“单点概率为零”预留直观基础。

\section{例题、习题与学习检查}
规划生日碰撞、无放回抽样、随机到达的会面问题三个例题，逐项写明试验机制、样本空间与目标事件。
证明题训练从公理推导结论；辨析题要求构造“各类结果不等可能”的反例。
完成本章后，应能说明一个概率答案所依赖的模型，而不只给出计数结果。






